\begin{abstract}
Classical traffic analysis explains how encrypted communication reveals activity. Which conditions carry over when an Agent forwards tool results to a model and executes privacy decoys as application work? We connect twenty-two mechanism anchors to fifteen direct Agent studies and forty-four settings through a transfer matrix, deployment access map and evidence-availability summary. The synthesis identifies two execution boundaries for protection. Feedback crosses separately controlled services, so a tool-side intervention must be evaluated against the model and recipient views that remain. Executed cover crosses into application work, so its privacy result needs an original-task constraint as well as a traffic budget. These extend established evaluation principles to the Agent execution path. Among four query/workflow privacy settings from two sources, protected-condition attack adaptation is located in two and matched original-task utility in one; none documents both in the retained records. The executed-decoy utility gap is supported by a published study; several model-view and adaptation measurements depend on preprints. A controlled feedback pilot supplies phase and complete-execution comparisons, with no multiplicity-adjusted primary rejection. The resulting design priorities are protection across services and privacy--utility measurement under the same intervention.
\end{abstract}
